Let us finally prove (i)\(\Rightarrow\)(v).  If \(Q\) is the radical of
\(G\), then
\(\dim(T)-\dim(Q)=1\), and \(Q\) is central in \(G\).  Hence
\(G=\underline{\operatorname{Centr}}_{G}(Q)\), which proves that \(Q\)
is a singular torus.  By \emph{Bible}, §11.3, Theorem 2,
\[
  W_G(T)=(\mathbb Z/2\mathbb Z)_k,
\]
and, by \emph{Bible}, §12.1, Lemma 1,
\(\dim(G)-\dim(T)=2\).  Thus \(G\) has at most two roots relative to
\(T\), while it has at least two, opposite to one another, by \SGARefLink{sga3:xix:paragraph:1-10}{1.10}.

\par\medskip\noindent
\SGAInlineRefTarget{sga3:xix:proposition:1-12}\textbf{\textit{\SGARefLink{sga3:xix:proposition:1:12}{Proposition 1.12}.}}---
\textit{Let \(k\) be an algebraically closed field, \(G\) a smooth
connected \(k\)-group, \(T\) a torus of \(G\), and \(R\) the set of
roots of \(G\) relative to \(T\).  Write
\[
  \mathfrak g
  =
  \mathfrak g^{0}
  \oplus
  \bigoplus_{\alpha\in R}\mathfrak g^{\alpha},
  \qquad
  \mathfrak g^{\alpha}\ne0,
\]
for the decomposition of \(\mathfrak g\) under
\(\operatorname{Ad}(T)\).  For every \(\alpha\in R\), put
\[
  T_{\alpha}=\operatorname{Ker}(\alpha)^{0},
  \qquad
  Z_{\alpha}=\underline{\operatorname{Centr}}_{G}(T_{\alpha}).
\]
\footnote{Editor's note: thus \(T_{\alpha}\) is the identity component
of \(\operatorname{Ker}(\alpha)\).}
The following conditions are equivalent:
\begin{enumerate}
\item[\SGAInlineRefTarget{sga3:xix:lemma:1:11:item:i}(i)] \(G\) is affine and reductive, and \(T\) is maximal.
\item[\SGAInlineRefTarget{sga3:xix:lemma:1:11:item:ii}(ii)] \(\mathfrak g^{0}=\mathfrak t\), and every \(Z_{\alpha}\),
  for \(\alpha\in R\), is reductive.
\item[\SGAInlineRefTarget{sga3:xix:lemma:1:11:item:iii}(iii)] \(\mathfrak g^{0}=\mathfrak t\); every
  \(\mathfrak g^{\alpha}\), for \(\alpha\in R\), has dimension \(1\);
  if \(\alpha,\beta\in R\) and \(\beta=q\alpha\) for some
  \(q\in\mathbb Q\), then \(q=\pm1\); and, for every \(\alpha\in R\),
  there exists \(w_{\alpha}\in G(k)\) which normalizes \(T\),
  centralizes \(T_{\alpha}\), but does not centralize \(T\).
\end{enumerate}
Moreover, under these conditions every \(Z_{\alpha}\) has semisimple
rank \(1\), and
\[
  \operatorname{Lie}(Z_{\alpha})
  =
  \mathfrak t\oplus
  \mathfrak g^{\alpha}\oplus
  \mathfrak g^{-\alpha}.
\]}
\SGARefTarget{sga3:xix:proposition:1:12}

\SGAProof{Proof}
Suppose first that \(G\) is affine and reductive and that \(T\) is
maximal.  Every \(Z_{\alpha}\) is then affine and reductive by \SGARefLink{sga3:xix:lemma:1:6:2:item:i}{1.6.2(i)},
with maximal torus \(T\); furthermore, \(T\) is its own centralizer by
\SGARefLink{sga3:xix:lemma:1:6:2:item:ii}{1.6.2(ii)}.  Thus \(\mathfrak g^{0}=\mathfrak t\), proving
\SGARefLink{sga3:xix:proposition:1:12:item:i}{(i)}\(\Rightarrow\)\SGARefLink{sga3:xix:proposition:1:12:item:ii}{(ii)}.

Conversely, the equality \(\mathfrak g^{0}=\mathfrak t\) implies that
\(T\) is maximal and that \(G\) is affine, by \SGARefLink{sga3:xix:paragraph:1-10}{1.10}.  Hence every
\(Z_{\alpha}\) is affine, smooth, and connected, by \SGARefLink{sga3:xix:paragraph:1-3}{1.3}.  In all cases,
\SGARefLink{sga3:ii:lemma:5-2-2}{Exposé II, 5.2.2} gives
\[
  \operatorname{Lie}(Z_{\alpha})
  =
  \mathfrak g^{T_{\alpha}}
  =
  \mathfrak g^{0}
  \oplus
  \bigoplus_{\beta\in R\cap\mathbb Q\alpha}
    \mathfrak g^{\beta}.
\]
Consequently,
\[
  \operatorname{Lie}(Z_{\alpha})
  \supset
  \mathfrak t\oplus
  \mathfrak g^{\alpha}\oplus
  \mathfrak g^{-\alpha},
\]
so \(Z_{\alpha}\ne T\).  Since \(T_{\alpha}\) is a central subtorus of
\(Z_{\alpha}\), of codimension \(1\) in \(T\), \SGARefLink{sga3:xix:lemma:1-11}{Lemma 1.11} applied to
\(Z_{\alpha}\) proves the equivalence of \SGARefLink{sga3:xix:proposition:1:12:item:ii}{(ii)} and (iii), together with
the additional assertions.

It remains to prove \SGARefLink{sga3:xix:proposition:1:12:item:ii}{(ii)}\(\Rightarrow\)\SGARefLink{sga3:xix:proposition:1:12:item:i}{(i)}.  We already know that \SGARefLink{sga3:xix:proposition:1:12:item:ii}{(ii)}
implies that \(T\) is maximal and that \(G\) is affine.  Let \(U\) be
the unipotent radical of \(G\).  It is invariant in \(G\), and its Lie
algebra is invariant under \(\operatorname{Ad}(T)\); hence
\[
  \operatorname{Lie}(U)
  =
  \bigl(\operatorname{Lie}(U)\cap\mathfrak g^{0}\bigr)
  \oplus
  \bigoplus_{\alpha\in R}
    \bigl(\operatorname{Lie}(U)\cap\mathfrak g^{\alpha}\bigr).
\]
By (\SGARefLink{sga3:xix:eq:1:3:1}{1.3.1}),
\[
\begin{aligned}
  U\cap T
    &=
    U\cap\underline{\operatorname{Centr}}_{G}(T)
     =
    \operatorname{rad}^{u}
      \bigl(\underline{\operatorname{Centr}}_{G}(T)\bigr)
     =
    \operatorname{rad}^{u}(T)
     =
    \{e\},\\
  U\cap Z_{\alpha}
    &=
    U\cap\underline{\operatorname{Centr}}_{G}(T_{\alpha})
     =
    \operatorname{rad}^{u}
      \bigl(\underline{\operatorname{Centr}}_{G}(T_{\alpha})\bigr)
     =
    \operatorname{rad}^{u}(Z_{\alpha})
     =
    \{e\}.
\end{aligned}
\]
It follows\footnote{Editor's note: if \(U\) and \(L\) are two subgroup
schemes of \(G\), then
\(\operatorname{Lie}(U)\cap\operatorname{Lie}(L)
=\operatorname{Lie}(U\cap L)\).} that
\[
\begin{aligned}
  \operatorname{Lie}(U)\cap\mathfrak g^{0}
    &=
    \operatorname{Lie}(U)\cap\mathfrak t
     =
    \operatorname{Lie}(U\cap T)
     =
    0,\\
  \operatorname{Lie}(U)\cap\mathfrak g^{\alpha}
    &\subset
    \operatorname{Lie}(U)\cap\operatorname{Lie}(Z_{\alpha})
     =
    \operatorname{Lie}(U\cap Z_{\alpha})
     =
    0.
\end{aligned}
\]
Thus \(\operatorname{Lie}(U)=0\), and therefore \(U=\{e\}\); in other
words, \(G\) is reductive.

\par\medskip\noindent
\SGAInlineRefTarget{sga3:xix:corollary:1-13}\textbf{\textit{\SGARefLink{sga3:xix:corollary:1:13}{Corollary 1.13}.}}---
\textit{Let \(k\), \(G\), \(T\), \(R\), and
\[
  \mathfrak g
  =
  \mathfrak g^{0}
  \oplus
  \bigoplus_{\alpha\in R}\mathfrak g^{\alpha}
\]
be as in \SGARefLink{sga3:xix:proposition:1-12}{Proposition 1.12}.  For every \(\alpha\in R\), let
\(T_{\alpha}\) be the maximal torus of \(\operatorname{Ker}(\alpha)\)
and put
\(Z_{\alpha}=\underline{\operatorname{Centr}}_{G}(T_{\alpha})\).
The following conditions are equivalent:
\begin{enumerate}
\item[\SGAInlineRefTarget{sga3:xix:proposition:1:12:item:i}(i)] \(G\) is affine and semisimple, and \(T\) is maximal.
\item[\SGAInlineRefTarget{sga3:xix:proposition:1:12:item:ii}(ii)] \(\mathfrak g^{0}=\mathfrak t\), every \(Z_{\alpha}\) is
  reductive, and
  \(\bigcap_{\alpha\in R}\operatorname{Ker}(\alpha)\) is finite.
\end{enumerate}
\footnote{Editor's note: this follows from the equalities
\(\operatorname{rad}(G)
=\underline{\operatorname{Centr}}(G)^{0}\) and
\(\underline{\operatorname{Centr}}(G)
=\bigcap_{\alpha\in R}\operatorname{Ker}(\alpha)\), established in
\SGARefLink{sga3:xix:lemma:1:6:2:item:iv}{1.6.2(iv)} and (\SGARefLink{sga3:xix:equation:1:10:3}{1.10.3}).}}
\SGARefTarget{sga3:xix:corollary:1:13}
